\documentclass[12pt,letterpaper]{article}

% Packages
\usepackage[utf8]{inputenc}
\usepackage[T1]{fontenc}
\usepackage{geometry}
\geometry{margin=1in}
\usepackage{graphicx}
\usepackage{booktabs}
\usepackage{longtable}
\usepackage{amsmath}
\usepackage{amssymb}
\usepackage{natbib}
\usepackage{hyperref}
\usepackage{xcolor}
\usepackage{lineno}
\usepackage{setspace}
\usepackage{caption}
\usepackage{subcaption}
\usepackage{siunitx}

% Line numbers for review
\linenumbers
\doublespacing

% Custom commands
\newcommand{\TODO}[1]{\textcolor{red}{[TODO: #1]}}
\newcommand{\CITE}{\textcolor{blue}{[CITE]}}

% Title
\title{Metal Contamination in Wildfire Ash: Geochemical Characterization and Remote Sensing Proxies from the 2025 Eaton Fire, Los Angeles County, California}

\author{
Isaac D. Medina\textsuperscript{1,*} \\
\textsuperscript{1}Department of Geography, University of California, Los Angeles \\
\textsuperscript{*}Corresponding author: imedina@ucla.edu
}

\date{Draft: December 2025}

\begin{document}

\maketitle

%==============================================================================
\begin{abstract}
%==============================================================================

Wildfire ash from burned structures in the wildland-urban interface (WUI) poses significant contamination risks due to the concentration of heavy metals from building materials, household products, and vehicles. The January 2025 Eaton Fire in Los Angeles County destroyed over 9,000 structures in the foothill communities of Altadena and Pasadena, generating widespread ash deposits overlying the Raymond Groundwater Basin. We characterized metal contamination in 39 ash samples using ICP-MS and portable XRF, evaluated spectral proxies from field spectroscopy (ASD), and integrated open-source geospatial data to identify contamination predictors. Lead concentrations ranged from 14 to 18,528 ppm (median: 61 ppm), with 15\% of samples exceeding the EPA residential soil screening level of 400 ppm. Portable XRF systematically overestimated Pb concentrations, requiring correction: $\text{Pb}_{\text{ICP-MS}} = (\text{Pb}_{\text{XRF}} - 104) / 3.9$ (R² = 0.26 after outlier removal). Residual analysis revealed that char content (spectral char index: $\rho$ = -0.63, p = 0.009) and sample brightness (NIR albedo: $\rho$ = +0.70, p = 0.003) significantly influenced XRF bias. Among open-source predictors, XRF-derived Fe\textsubscript{2}O\textsubscript{3}\% showed the strongest correlation with Pb ($\rho$ = -0.63, p < 0.001), while property density was positively correlated ($\rho$ = +0.51, p = 0.001). Contrary to expectations, pre-1978 housing age was not a significant predictor of Pb contamination in this study area. These findings demonstrate that (1) XRF field screening requires matrix-specific correction factors, (2) spectral properties can inform XRF bias estimation, and (3) open-source geospatial data provide useful but imperfect proxies for contamination risk assessment.

\textbf{Keywords:} wildfire ash, lead contamination, XRF correction, spectral analysis, Raymond Basin, Eaton Fire

\end{abstract}

%==============================================================================
\section{Introduction}
%==============================================================================

\subsection{Wildfire-Urban Interface Contamination}

The increasing frequency and intensity of wildfires in California's wildland-urban interface (WUI) has created a growing environmental health crisis \citep{Radeloff2018}. When fires burn through residential and commercial areas, the combustion of building materials, household products, vehicles, and landscaping generates ash deposits containing elevated concentrations of heavy metals, asbestos, and organic contaminants \citep{Plumlee2007, Quarles2020}. Unlike wildland fire ash, which primarily reflects the geochemistry of burned vegetation and soils, structure fire ash can contain lead (Pb) from paint, pipes, and solder; arsenic (As) from treated wood; copper (Cu) and zinc (Zn) from electrical wiring and galvanized metals; cadmium (Cd) from batteries and plastics; and chromium (Cr) from chrome plating and pressure-treated lumber \citep{Neitzel2021}.

The January 2025 Eaton Fire exemplifies the scale of this emerging hazard. Igniting in Eaton Canyon on January 7, 2025, during an extreme Santa Ana wind event, the fire rapidly burned through the foothill communities of Altadena and portions of Pasadena, ultimately destroying over 9,000 structures across approximately 14,000 acres (5,670 ha) before full containment on January 31, 2025 \CITE. The fire caused 19 confirmed fatalities and became the fifth deadliest and second most destructive wildfire in California history \CITE. The widespread distribution of potentially contaminated ash across residential neighborhoods, parks, and recharge zones for the Raymond Groundwater Basin created an urgent need for rapid assessment of contamination levels and spatial patterns.

Current approaches to post-fire contamination assessment rely primarily on laboratory analysis of soil and ash samples, which is accurate but time-consuming and expensive. Portable X-ray fluorescence (XRF) instruments offer a rapid field screening alternative but are known to exhibit matrix-dependent biases in heterogeneous materials such as wildfire ash \citep{Kalnicky2001}. Understanding and correcting these biases is essential for effective rapid response.

\subsection{Study Objectives}

This study addresses four primary objectives:

\begin{enumerate}
    \item Characterize metal contamination levels in Eaton Fire ash and compare to regulatory thresholds
    \item Evaluate portable XRF as a field screening tool and develop empirical correction factors
    \item Identify spectral and open-source geospatial data proxies for contamination prediction
    \item Establish the geogenic baseline for distinguishing anthropogenic enrichment
\end{enumerate}

%==============================================================================
\section{Study Area}
%==============================================================================

\subsection{Geological and Hydrogeological Setting}

The Eaton Fire study area lies within the Raymond Groundwater Basin in the San Gabriel Valley, at the base of the San Gabriel Mountains in Los Angeles County, California (Figure~\ref{fig:study_area}). Sample locations span 34.127°N to 34.207°N latitude and 118.194°W to 118.020°W longitude, encompassing the communities of Altadena and portions of Pasadena at elevations between approximately 360 and 550 m above sea level.

\subsubsection{Geological Framework}

The study area occupies the transition zone between the crystalline basement of the San Gabriel Mountains and the alluvial sediments of the Raymond Basin. The San Gabriel Mountains represent a large fault-bounded block composed primarily of Precambrian gneiss and Mesozoic granitic rocks, uplifted along the Sierra Madre-Cucamonga Fault Zone during Miocene through Quaternary tectonism \citep{Yerkes1985}. The steep, deeply dissected terrain of the range---with slopes exceeding 70\% in Eaton Canyon---efficiently channels storm runoff and sediment transport to the basin below.

The foothill communities of Altadena and Pasadena are built upon coalescing alluvial fan deposits sourced from major drainages including Eaton Wash, the Arroyo Seco, and numerous smaller canyons. These deposits comprise:

\begin{itemize}
    \item \textbf{Holocene alluvium}: Unconsolidated sand, gravel, and cobbles in active channels; up to 20 m thick
    \item \textbf{Upper Pleistocene deposits}: Semi-consolidated alluvial fan sediments; gravelly sands with interbedded silts
    \item \textbf{Older alluvium}: Moderately consolidated fan deposits extending to depths of 150+ m
\end{itemize}

The Raymond Fault, an east-northeast trending left-lateral structure with minor reverse slip, traverses the southern portion of the study area. This fault acts as a partial groundwater barrier, separating the unconfined Raymond Basin to the north from the semi-confined aquifer systems of the Central Basin to the south \citep{DWR2003}.

\subsubsection{Hydrogeological Setting}

The Raymond Basin is a 104 km² (40 mi²) groundwater basin underlying Pasadena, Altadena, Sierra Madre, and portions of Arcadia (Table~\ref{tab:hydrogeology}). The basin is bounded to the north by the San Gabriel Mountain front (granitic bedrock), to the east by less permeable older sediments, to the south by the Raymond Fault, and to the west by the Arroyo Seco \citep{DWR2003}.

The aquifer system comprises unconsolidated to semi-consolidated alluvial deposits with an average thickness of approximately 170 m (550 ft) and maximum thickness exceeding 300 m near the mountain front. Porosity ranges from 3--35\% with an average of approximately 15\%. Historical storage capacity has been estimated at approximately 1,000,000 acre-feet \citep{DWR1961}.

Groundwater recharge occurs primarily via infiltration from the Arroyo Seco and Eaton Wash spreading basins, supplemented by direct precipitation infiltration. Average annual precipitation in the basin is approximately 380 mm (15 in). Groundwater in the Raymond Basin is characteristically calcium-bicarbonate type, reflecting dissolution of calcium-bearing minerals in the alluvial sediments. Natural fluoride concentrations occasionally exceed 1.6 mg/L near the mountain front (maximum 3.1 mg/L), indicating geogenic enrichment from granitic source rocks \citep{DWR2003}.

The Raymond Basin was adjudicated in 1944 (\textit{City of Pasadena v. City of Alhambra})---the first groundwater adjudication in California---and currently supplies approximately one-third of Pasadena's municipal water supply \citep{RBMB2020}. Notably, the NASA Jet Propulsion Laboratory (JPL), located within the study area, has historically contributed perchlorate and volatile organic compound (VOC) contamination to groundwater in the Monk Hill subarea, necessitating ongoing treatment at the Monk Hill Water Facility \citep{USEPA2018}.

\begin{table}[htbp]
\centering
\caption{Hydrogeological characteristics of the Raymond Basin}
\label{tab:hydrogeology}
\begin{tabular}{lll}
\toprule
\textbf{Parameter} & \textbf{Value} & \textbf{Source} \\
\midrule
Basin area & 104 km² (40 mi²) & DWR 2003 \\
Average thickness & 170 m (550 ft) & DWR 2003 \\
Maximum thickness & >300 m & DWR 2003 \\
Storage capacity & 1,000,000 acre-ft & DWR 1961 \\
Average porosity & 15\% (range 3--35\%) & DWR 2003 \\
Annual precipitation & 380 mm (15 in) & DWR 2004 \\
Groundwater type & Ca-HCO\textsubscript{3} & DWR 2003 \\
Adjudication year & 1944 & RBMB 2020 \\
\bottomrule
\end{tabular}
\end{table}

\subsubsection{Geogenic Metal Background}

Understanding the natural geochemical background is critical for distinguishing anthropogenic contamination in wildfire ash. The granitic and gneissic source rocks of the San Gabriel Mountains typically contain lead (Pb) at 10--30 ppm (crustal average $\sim$14 ppm), copper (Cu) at 10--50 ppm, zinc (Zn) at 40--80 ppm, and arsenic (As) at <5 ppm \citep{Taylor1964, Reimann2018}. Alluvial soils derived from these sources are expected to have similar or slightly lower concentrations due to weathering and dilution.

The Ca/Mg ratio in local sediments (typically 2--6) provides a geogenic signature useful for distinguishing natural from anthropogenic sources. Similarly, the Fe/Al ratio reflects the mineralogy of source rocks and can indicate the presence of iron-rich anthropogenic materials when elevated above background levels.

\subsubsection{Anthropogenic Legacy}

The study area has been continuously developed since the late 19th century, with most residential construction occurring between 1920 and 1970. This period corresponds to peak use of lead-based paints (phased out 1978), leaded gasoline emissions (phased out 1996), and lead water service lines (common pre-1950) \citep{Mielke1999}.

The broader Los Angeles Basin has documented legacy lead contamination from multiple industrial sources. Most notably, the Exide Technologies battery recycling facility in Vernon---approximately 20 km south of the study area---contaminated over 4,800 properties requiring soil remediation, with atmospheric deposition affecting neighborhoods throughout southeastern Los Angeles County \citep{DTSC2020}. While the Eaton Fire study area lies north of the primary Exide impact zone, regional atmospheric deposition patterns may contribute to elevated baseline lead levels in urban soils throughout the region, complicating the assessment of fire-specific contamination.

\subsection{Fire Characteristics}

The Eaton Fire ignited on January 7, 2025, in Eaton Canyon during an extreme Santa Ana wind event characterized by sustained winds of 80--100 km/h with gusts exceeding 150 km/h. Key fire characteristics include:

\begin{itemize}
    \item \textbf{Duration}: 24 days (contained January 31, 2025)
    \item \textbf{Area burned}: $\sim$14,000 acres (5,670 ha)
    \item \textbf{Structures destroyed}: >9,000
    \item \textbf{Fatalities}: 19 confirmed
    \item \textbf{Burn severity}: Predominantly moderate to high in WUI areas
\end{itemize}

The fire burned through established residential neighborhoods with homes ranging from 1890s Victorian-era construction to modern builds, representing a diversity of building materials and potential contaminant sources. The extreme fire behavior---driven by dry fuels, low humidity, and high winds---resulted in complete combustion of many structures, concentrating non-volatile metals in the resulting ash.

%==============================================================================
\section{Methods}
%==============================================================================

\subsection{Sample Collection}

A total of 39 ash samples and 3 reference soil samples were collected from the Eaton Fire burn area between January 15 and February 8, 2025. Sample locations were selected to represent the range of burn severities and structure types across the affected area. At each site, GPS coordinates were recorded and a composite sample was collected from the ash layer overlying the structure footprint, avoiding visible debris (e.g., metal fragments, glass) that could bias bulk geochemical analysis.

Samples were collected using stainless steel tools, placed in pre-cleaned polyethylene bags, and stored at ambient temperature prior to analysis. Field observations including structure type, apparent burn severity, and presence of specific materials (e.g., vehicles, outbuildings) were recorded for each site.

\subsection{Laboratory Analysis}

\subsubsection{ICP-MS Analysis}

Samples were digested using EPA Method 3051A (microwave-assisted acid digestion) and analyzed by inductively coupled plasma mass spectrometry (ICP-MS) for a suite of 50+ elements including Pb, As, Cd, Cu, Zn, Cr, Ni, and rare earth elements (REEs). Detection limits were typically <0.1 ppm for most trace metals. Quality assurance included analysis of certified reference materials (CRMs), method blanks, and duplicate samples at a frequency of 10\%.

\subsubsection{Portable XRF Analysis}

Field XRF measurements were made using a \TODO{instrument model} portable XRF analyzer operated in soil mode. Major element oxides (SiO\textsubscript{2}, Al\textsubscript{2}O\textsubscript{3}, Fe\textsubscript{2}O\textsubscript{3}, CaO, MgO, K\textsubscript{2}O, TiO\textsubscript{2}, etc.) and trace metals (Pb, Zn, Cu) were measured with typical integration times of 60 seconds per sample. XRF results for major oxides are reported in weight percent (wt\%), while trace metals reported by the instrument were converted from wt\% to ppm by multiplying by 10,000 for comparison with ICP-MS results.

\subsection{Spectroscopic Measurements}

\subsubsection{ASD Field Spectroscopy}

Visible to shortwave infrared (VNIR-SWIR) reflectance spectra were measured for a subset of samples (n = 21) using an ASD FieldSpec 4 spectroradiometer (350--2500 nm). Spectra were collected under controlled illumination conditions using a contact probe. Multiple measurements (5--10) were averaged for each sample to reduce noise.

Spectral indices were calculated from the reflectance data including:
\begin{itemize}
    \item \textbf{Fe index}: Ratio of reflectance at 870 nm to 650 nm, sensitive to iron oxide content
    \item \textbf{Char index}: Overall reflectance darkness (1 - mean albedo), indicating carbonaceous material
    \item \textbf{Carbonate index}: Depth of 2.3 \textmu m absorption feature
    \item \textbf{Spectral slopes}: Slope of reflectance in VIS (450--700 nm) and NIR (700--1000 nm) regions
\end{itemize}

\subsubsection{FTIR Analysis}

Fourier transform infrared (FTIR) spectra were measured on dried, ground samples to characterize organic functional groups and mineral phases. Spectra were collected in attenuated total reflectance (ATR) mode from 4000--400 cm\textsuperscript{-1}.

\subsection{Open-Source Data Integration}

\subsubsection{Property Age Data}

Property construction dates were obtained from the Los Angeles County Assessor's database via the LA County Parcel API. For each sample location, parcels within a 250 m radius were queried and the percentage of structures built before 1978 (pre-lead paint ban era) was calculated. Mean construction year and total property count (as a proxy for housing density) were also recorded.

\subsubsection{Census Data}

Census tract identifiers were obtained for each sample location using the U.S. Census Bureau Geocoder API. Tract-level demographic and housing characteristics were extracted from the American Community Survey for potential correlation with contamination patterns.

\subsubsection{Geogenic Ratios}

Geogenic indicator ratios were calculated from XRF major element data:
\begin{itemize}
    \item Ca/Mg ratio: Indicator of calcium carbonate vs. dolomite influence
    \item Fe/Al ratio: Indicator of iron enrichment relative to aluminosilicates
    \item Si/Al ratio: Indicator of quartz content relative to clay minerals
\end{itemize}

\subsection{Statistical Analysis}

All statistical analyses were performed in Python using scipy, numpy, and scikit-learn libraries. Spearman rank correlations were used to assess relationships between metal concentrations and potential predictor variables, as metal distributions were non-normal. Linear regression was used to develop XRF correction models, with residual analysis to identify matrix effects. Significance was assessed at $\alpha$ = 0.05 with Bonferroni correction applied for multiple comparisons where appropriate.

%==============================================================================
\section{Results}
%==============================================================================

\subsection{Metal Contamination Levels}

Lead (Pb) concentrations in ash samples ranged from 14 to 18,528 ppm with a median of 61 ppm and mean of 638 ppm (Table~\ref{tab:metal_summary}). The distribution was strongly right-skewed, with a small number of highly contaminated samples driving the elevated mean. Fifteen percent of samples exceeded the EPA residential soil screening level (RSL) of 400 ppm for Pb.

\begin{table}[htbp]
\centering
\caption{Summary statistics for metal concentrations in Eaton Fire ash (n=39)}
\label{tab:metal_summary}
\begin{tabular}{lrrrrrr}
\toprule
\textbf{Metal} & \textbf{Min} & \textbf{Median} & \textbf{Mean} & \textbf{Max} & \textbf{RSL} & \textbf{\% >RSL} \\
\midrule
Pb (ppm) & 14 & 61 & 638 & 18,528 & 400 & 15 \\
As (ppm) & 3.0 & 6.1 & 6.5 & 12.2 & 0.68 & 100 \\
Cu (ppm) & 14 & 61 & 207 & 3,650 & 3,100 & 3 \\
Zn (ppm) & 100 & 333 & 1,310 & 23,829 & 23,000 & 3 \\
Cd (ppm) & 0.1 & 0.3 & 0.5 & 1.5 & 70 & 0 \\
Cr (ppm) & 22 & 43 & 55 & 332 & 120,000 & 0 \\
\bottomrule
\end{tabular}
\end{table}

\subsection{XRF-ICP-MS Correction Model}

Comparison of portable XRF and ICP-MS measurements revealed systematic overestimation of Pb by XRF (Figure~\ref{fig:xrf_correction}). Using all data (n = 37), the linear relationship was:

\begin{equation}
\text{Pb}_{\text{XRF}} = 1.71 \times \text{Pb}_{\text{ICP-MS}} + 337 \quad (R^2 = 0.98)
\end{equation}

However, this model was heavily influenced by extreme outliers. After removing 7 statistical outliers identified by the IQR method (samples with Pb > 134 ppm by ICP-MS or XRF > 805 ppm), the relationship for typical samples (n = 30) was:

\begin{equation}
\text{Pb}_{\text{XRF}} = 3.90 \times \text{Pb}_{\text{ICP-MS}} + 104 \quad (R^2 = 0.26)
\end{equation}

The substantially lower R² for the cleaned dataset indicates high variability in XRF measurements for typical ash samples, likely due to matrix effects. The practical correction formula for converting XRF field measurements to ICP-MS equivalent values is:

\begin{equation}
\text{Pb}_{\text{ICP-MS}} \approx \frac{\text{Pb}_{\text{XRF}} - 104}{3.9}
\end{equation}

For field screening purposes, the XRF threshold corresponding to the EPA RSL of 400 ppm is:

\begin{equation}
\text{Pb}_{\text{XRF, threshold}} = 3.90 \times 400 + 104 = 1,662 \text{ ppm}
\end{equation}

\subsection{Residual Analysis and Matrix Effects}

Analysis of residuals from the XRF-ICP-MS regression revealed significant correlations with both spectral and matrix properties (Table~\ref{tab:residuals}). Among spectral features, NIR albedo showed the strongest correlation ($\rho$ = +0.70, p = 0.003), followed by visible spectral slope ($\rho$ = +0.67, p = 0.005) and char index ($\rho$ = -0.63, p = 0.009). The negative correlation with char index indicates that darker, more char-rich samples tend to show lower XRF values relative to ICP-MS (i.e., underestimation), while brighter samples show higher XRF values (overestimation).

Among matrix properties derived from XRF major elements, MgO content showed a significant negative correlation with residuals ($\rho$ = -0.53, p = 0.014), suggesting that magnesium-rich matrices attenuate X-ray signals.

\begin{table}[htbp]
\centering
\caption{Spearman correlations between XRF residuals and predictor variables (n = 16--21)}
\label{tab:residuals}
\begin{tabular}{llrrr}
\toprule
\textbf{Type} & \textbf{Variable} & \textbf{$\rho$} & \textbf{p-value} & \textbf{Significance} \\
\midrule
Spectral & Albedo\_NIR & +0.70 & 0.003 & ** \\
Spectral & Slope\_VIS & +0.67 & 0.004 & ** \\
Spectral & Char\_index & -0.63 & 0.009 & ** \\
Spectral & Albedo\_VIS & +0.30 & 0.264 & \\
Matrix & MgO\% & -0.58 & 0.018 & * \\
Matrix & Fe\textsubscript{2}O\textsubscript{3}\% & -0.07 & 0.795 & \\
Matrix & CaO\% & -0.16 & 0.564 & \\
\bottomrule
\end{tabular}
\end{table}

\subsection{Open-Source Predictors of Contamination}

Among open-source data sources evaluated as potential predictors of Pb contamination, XRF-derived Fe\textsubscript{2}O\textsubscript{3}\% showed the strongest correlation ($\rho$ = -0.63, p < 0.001), indicating that lower iron content is associated with higher Pb (Figure~\ref{fig:predictors}). Property density (count within 250 m) was positively correlated with Pb ($\rho$ = +0.51, p = 0.001), consistent with the expectation that denser development corresponds to more potential contamination sources.

Notably, the percentage of pre-1978 housing---hypothesized to predict lead contamination due to lead-based paint---was \textit{not} significantly correlated with Pb ($\rho$ = -0.12, p = 0.49). However, pre-1978 housing percentage was significantly correlated with nickel (Ni) concentrations ($\rho$ = +0.48, p = 0.003), suggesting different source attribution for these metals.

The geogenic Ca/Mg ratio showed moderate positive correlation with Zn ($\rho$ = +0.49, p = 0.002) but not with Pb, suggesting potential geochemical controls on Zn distribution independent of anthropogenic sources.

%==============================================================================
\section{Discussion}
%==============================================================================

\subsection{Contamination Sources and Patterns}

The wide range of Pb concentrations observed (14--18,528 ppm) reflects the heterogeneous nature of structure fire ash, where localized ``hot spots'' associated with specific materials (e.g., painted surfaces, batteries, lead pipes) occur within a matrix of lower-contamination ash. The median concentration of 61 ppm is below the EPA RSL of 400 ppm but exceeds typical geogenic background levels of 10--30 ppm, indicating widespread low-level anthropogenic enrichment even in samples not classified as hazardous.

The lack of correlation between Pb and pre-1978 housing age is surprising given the well-documented association between older housing and lead contamination in urban soils \citep{Mielke1999}. Several factors may explain this finding:

\begin{enumerate}
    \item \textbf{Combustion homogenization}: Fire may redistribute and mix materials from different construction eras within a neighborhood
    \item \textbf{Non-paint sources}: Lead in structure fire ash may derive substantially from non-paint sources (pipes, solder, electronics) not strictly correlated with construction date
    \item \textbf{Renovation history}: Older homes may have been renovated with lead-free materials, while some newer homes may contain lead from imported goods
    \item \textbf{Spatial confounding}: Housing age may be correlated with other factors (e.g., lot size, socioeconomic status) that influence contamination patterns
\end{enumerate}

The significant correlation between housing density and Pb supports a ``mass of materials'' interpretation, where contamination scales with the total amount of combustible materials regardless of specific sources.

\subsection{XRF Limitations and Recommendations}

The high variability in XRF measurements for typical ash samples (R² = 0.26 after outlier removal) presents challenges for field screening applications. The systematic overestimation (slope $\sim$3.9) combined with a positive intercept ($\sim$104 ppm) means that XRF will classify some samples as potentially hazardous that would test below the RSL by laboratory methods.

For conservative screening, this positive bias may be acceptable, as it errs on the side of caution. However, the dependence of bias on sample properties---particularly char content---introduces systematic errors that could lead to inconsistent classification. Darker, char-rich samples will tend to be under-flagged, while brighter samples will be over-flagged.

We recommend the following protocol for XRF field screening of wildfire ash:

\begin{enumerate}
    \item Apply the empirical correction: $\text{Pb}_{\text{est}} = (\text{Pb}_{\text{XRF}} - 104) / 3.9$
    \item Use an XRF threshold of 1,662 ppm as the primary screening criterion
    \item Visually assess sample brightness and adjust interpretation accordingly:
    \begin{itemize}
        \item Dark/char-rich samples: True contamination may be higher than XRF indicates
        \item Bright/ash-white samples: True contamination likely lower than XRF indicates
    \end{itemize}
    \item Confirm borderline samples (XRF 1,200--2,500 ppm) with laboratory analysis
\end{enumerate}

\subsection{Remote Sensing Potential}

The significant correlations between spectral properties and XRF bias suggest potential for remote sensing-informed correction of field measurements. Imaging spectrometers such as AVIRIS-NG can map char content and surface brightness at meter-scale resolution over large areas \citep{Green1998}. Integration of airborne spectral data with ground-based XRF screening could improve contamination mapping by:

\begin{enumerate}
    \item Providing spatially continuous estimates of matrix properties affecting XRF accuracy
    \item Identifying high-char zones where XRF underestimation is expected
    \item Enabling stratified sampling designs that account for spectral variability
\end{enumerate}

\subsection{Groundwater Vulnerability}

The location of the Eaton Fire burn area over the Raymond Groundwater Basin raises concerns about potential metal transport to groundwater through post-fire runoff and infiltration. The spreading basins along Eaton Wash and the Arroyo Seco that provide primary recharge to the basin could receive contaminated sediment during winter storms. Given that the Raymond Basin supplies approximately one-third of Pasadena's water supply, monitoring of recharge water quality is warranted.

However, several factors may limit groundwater impacts:
\begin{itemize}
    \item Lead is relatively immobile in soite environments at circumneutral pH
    \item The thick unsaturated zone (tens of meters) provides opportunity for attenuation
    \item Engineered debris basins may capture contaminated sediment before it reaches spreading grounds
\end{itemize}

Long-term monitoring of groundwater Pb concentrations downgradient of the burn area will be necessary to assess actual impacts.

%==============================================================================
\section{Conclusions}
%==============================================================================

Analysis of 39 ash samples from the 2025 Eaton Fire reveals the following key findings:

\begin{enumerate}
    \item \textbf{Contamination levels}: Pb concentrations ranged from 14 to 18,528 ppm (median 61 ppm), with 15\% exceeding the EPA RSL of 400 ppm. The distribution is highly skewed with localized hot spots.

    \item \textbf{XRF correction}: Portable XRF systematically overestimates Pb in ash samples. The recommended correction is: $\text{Pb}_{\text{ICP-MS}} = (\text{Pb}_{\text{XRF}} - 104) / 3.9$. For field screening, use an XRF threshold of 1,662 ppm to identify samples likely exceeding the 400 ppm RSL.

    \item \textbf{Matrix effects}: Char content introduces systematic bias in XRF measurements. Dark, char-rich samples show relative underestimation while bright samples show overestimation. These effects are detectable via spectral measurements (NIR albedo $\rho$ = +0.70, char index $\rho$ = -0.63).

    \item \textbf{Open-source predictors}: Fe\textsubscript{2}O\textsubscript{3}\% ($\rho$ = -0.63) and property density ($\rho$ = +0.51) are the strongest predictors of Pb contamination. Housing age (pre-1978 construction) is \textit{not} a significant predictor in this study area.

    \item \textbf{Geogenic baseline}: Background Pb from San Gabriel Mountains-derived alluvium is 10--30 ppm. Median ash Pb (61 ppm) indicates widespread low-level anthropogenic enrichment.
\end{enumerate}

These findings demonstrate the complexity of post-fire contamination assessment and highlight opportunities for integrating field spectroscopy and open-source geospatial data with traditional sampling approaches to improve rapid response capabilities.

%==============================================================================
% References
%==============================================================================

\bibliographystyle{agsm}
\bibliography{references}

%==============================================================================
% Figures
%==============================================================================

\clearpage
\section*{Figures}

\begin{figure}[htbp]
    \centering
    \fbox{\parbox{0.8\textwidth}{\centering\vspace{3cm}[Figure 1: Study area map]\vspace{3cm}}}
    \caption{Study area showing sample locations, Raymond Basin extent, Raymond Fault, and major drainages. Inset shows location within Los Angeles County.}
    \label{fig:study_area}
\end{figure}

\begin{figure}[htbp]
    \centering
    \fbox{\parbox{0.8\textwidth}{\centering\vspace{3cm}[Figure 2: XRF correction plots]\vspace{3cm}}}
    \caption{XRF vs ICP-MS Pb concentrations. (A) All data showing regression line and outliers. (B) Data with outliers removed showing higher scatter. (C) Residual distribution. (D) Corrected XRF values vs true ICP-MS values.}
    \label{fig:xrf_correction}
\end{figure}

\begin{figure}[htbp]
    \centering
    \fbox{\parbox{0.8\textwidth}{\centering\vspace{3cm}[Figure 3: Predictor correlations]\vspace{3cm}}}
    \caption{Correlations between Pb concentration and potential predictors. (A) Fe\textsubscript{2}O\textsubscript{3}\%. (B) Property density. (C) Pre-1978 housing percentage. (D) Correlation matrix heatmap.}
    \label{fig:predictors}
\end{figure}

\end{document}
